The pair sources, the bilinear $B$ and the first-order field equation are those of Theorem \claimref{synthesis.gravitational-pair-production-from-the-operator-source-at-first-order} (its clauses (i) to (iii)); the first-order curvature map and its kernels are those of Theorem \claimref{synthesis.weyl-flat-variations-carry-only-the-stress-trace-and-the-whole-stress-clause-is-the-trace-free-channel}. Throughout, $k$ is fixed with $Q(k) = s > 0$, $k^{\perp}$ is the space of vectors orthogonal to $k$, of dimension three and negative definite, and $\{e_{i}\}$, $i = 1, 2, 3$, is an orthonormal frame of it, $e^{i} = -e_{i}$. Every statement below is covariant under rotors of $\Cl(1,3)$ and homogeneous under rescaling of the momenta, and every pair $(p_{1}, p_{2})$ of null vectors with $Q(p_{1} + p_{2}) > 0$ is the image of $(e_{0} + e_{3}, e_{0} - e_{3})$ under a rotor and a rescaling; so a statement about one pair that is verified at this configuration holds for all.

\emph{(i).} A source is symmetric and kills $k$, so it maps into $k^{\perp}$ and is fixed by its restriction to $k^{\perp}$, a symmetric map of a three-dimensional space: dimension $6$. Since $\mathcal{T}(e_{i}) \in k^{\perp}$, $\ip{\mathcal{T}}{\mathcal{T}} = \sum_{i} \mathcal{T}(e_{i}) \cdot \mathcal{T}(e^{i}) = -\sum_{i} Q(\mathcal{T}(e_{i})) \ge 0$, with equality only for $\mathcal{T} = 0$. $P_{k}$ is the identity of $k^{\perp}$, so $\tr P_{k} = \ip{P_{k}}{P_{k}} = 3$ and $\ip{\mathcal{T}}{P_{k}} = \sum_{i} \mathcal{T}(e_{i}) \cdot e^{i} = \tr\mathcal{T}$. Hence $\mathcal{T}^{\circ}$ is trace-free and orthogonal to $P_{k}$, and with $t_{n} = \tr\mathcal{T}_{n}$,
$\ip{\mathcal{T}_{1}}{\mathcal{T}_{2}} = \ip{\mathcal{T}_{1}^{\circ}}{\mathcal{T}_{2}^{\circ}} + \tfrac19 t_{1}t_{2} \cdot 3$, so $B = \ip{\mathcal{T}_{1}^{\circ}}{\mathcal{T}_{2}^{\circ}} + (\tfrac13 - \tfrac12)\,t_{1}t_{2}$.

\emph{(ii).} The rotors fixing $k$ are the rotors of the even subalgebra of $\Cl(k^{\perp})$, a connected group, so an element of $V_{k}$ or of $V_{k} \otimes V_{k}$ is invariant under all of them exactly when it is annihilated by $L_{b}$ (acting as a derivation on the tensor square) for the three unit bivectors $b = e_{1}e_{2}, e_{2}e_{3}, e_{3}e_{1}$. This is a finite linear problem, solved exactly by the script at $k = 2e_{0}$: the common kernel of the three $L_{b}$ on $V_{k}$ is zero, and on $V_{k} \otimes V_{k}$, of dimension $25$, the invariance equations have rank $24$. The pairing is invariant ($\ip{L_{b}A}{C} + \ip{A}{L_{b}C} = 0$, checked on a basis), so its tensor $\sum_{\alpha} E_{\alpha} \otimes E_{\alpha}$, for an orthonormal basis $E_{\alpha}$ of $V_{k}$, spans the line. Now $M(C, C') = \langle R \cdot C \otimes R \cdot C' \rangle_{R}$ is invariant, so $M(C, C') = \lambda\sum_{\alpha} E_{\alpha} \otimes E_{\alpha}$; contracting the two factors with the pairing, and using its invariance, $\ip{C}{C'} = 5\lambda$. Pairing $M(C, C')$ with $A \otimes A'$ gives the first identity. The mean $\langle R \cdot C \rangle_{R}$ is an invariant element of $V_{k}$, so it is zero, which is the second.

\emph{(iii).} Scalar pair. For $a \in k^{\perp}$, $a \cdot p_{2} = -a \cdot p_{1}$ and $a \cdot p_{1} = \tfrac12\,a \cdot d$, so the stored source is $\mathcal{T}(a) = \tfrac12 (a \cdot d)\,d + \tfrac{s}{2}(1 - 4\xi)\,a$ on $k^{\perp}$, with $Q(d) = -2\,p_{1} \cdot p_{2} = -s$. It is built from $d$ and $P_{k}$, which $b$ leaves fixed ($b \times d = 0$), so it has weight zero. Its trace is $-\tfrac{s}{2} + \tfrac{3s}{2}(1 - 4\xi) = (1 - 6\xi)s$, and $\mathcal{T}^{\circ}(a) = \tfrac12\bigl((a \cdot d)\,d + \tfrac{s}{3}P_{k}(a)\bigr)$ has norm $\tfrac14\bigl(s^{2} - \tfrac23 s^{2} + \tfrac13 s^{2}\bigr) = s^{2}/6$.

Chiral pair. Write a real vector as $w = \alpha p_{1} + \beta p_{2} + w_{\perp}$ with $w_{\perp}$ in the plane of $b$. Then $\br{p_{2}wp_{1}w}_{0} = 2(p_{2} \cdot w)(p_{1} \cdot w) - (p_{1} \cdot p_{2})Q(w) = -(p_{1} \cdot p_{2})Q(w_{\perp})$, so the stored spin sum $2\br{p_{2}wp_{1}w}_{0}$ (half the sum of Theorem \claimref{synthesis.electron-positron-to-muon-pairs-in-the-spacetime-algebra-lsz-tree-cross-section-and-one-loop-unitarity}(1) at zero mass, as the stored theorem takes it) equals $-s\,Q(w_{\perp}) = (j_{1} \cdot w)^{2} + (j_{2} \cdot w)^{2}$ for any orthogonal pair $j_{1}, j_{2}$ of vectors in that plane with $Q(j_{1}) = Q(j_{2}) = -s$: these are the two real components of the pair's transition current, and $b \times j_{1} = j_{2}$, $b \times j_{2} = -j_{1}$ for the orientation of $b$ that is chosen so. Since $b \times d = 0$, $L_{b}\mathcal{T}_{1} = \mathcal{T}_{2}$ and $L_{b}\mathcal{T}_{2} = -\mathcal{T}_{1}$: a pair of weight one. $\tr\mathcal{T}_{n} = \tfrac12\,j_{n} \cdot d = 0$; $\ip{\mathcal{T}_{m}}{\mathcal{T}_{n}} = \tfrac{1}{16}\bigl(2\,(j_{m} \cdot j_{n})\,Q(d) + 2\,(j_{m} \cdot d)(j_{n} \cdot d)\bigr) = \tfrac18\,s^{2}\,\delta_{mn}$, so the norms sum to $s^{2}/4$.

Vector pair. The four sources are linear in each polarisation, and the polarisations of each quantum form a pair of weight one about $b$, so their span carries weights $0, 0$ and a pair of weight two. The script computes them at the reference configuration: the combinations $\epsilon_{1} = \epsilon_{2}$ summed over the frame and the antisymmetric one vanish identically, the two others satisfy $L_{b}L_{b}\mathcal{T} = -4\mathcal{T}$, every source is trace-free (the Maxwell stress of Fact \claimref{synthesis.imported-maxwell-stress-spacetime-algebra}) and conserved, each has norm $s^{2}/2$, and $\br{F_{1}F_{2}}_{0}$ is $s/2$ on equal polarisations and zero on unequal ones.

\emph{(iv).} Fix the first pair along a reference axis, with sources $A_{\alpha}$, and let the second pair point along a direction $n$. Its sources are $R \cdot C_{\beta}$ for any rotor $R$ fixing $k$ that carries the reference axis to $n$, the choice of $R$ changing them by a rotation about the axis, which by (iii) mixes the states orthogonally and leaves $\sum_{\beta}$ of any quadratic expression unchanged. So $\sum_{\alpha\beta} B(A_{\alpha}, R \cdot C_{\beta})^{2}$ depends on $R$ only through $n$, and its integral over the sphere is $4\pi$ times its mean over rotors. By (i), $B(A_{\alpha}, R \cdot C_{\beta}) = \ip{A_{\alpha}^{\circ}}{R \cdot C_{\beta}^{\circ}} - \tfrac16 t_{\alpha}t_{\beta}$; by (ii) the mean of its square is $\tfrac15\ip{A_{\alpha}^{\circ}}{A_{\alpha}^{\circ}}\ip{C_{\beta}^{\circ}}{C_{\beta}^{\circ}} + \tfrac{1}{36}t_{\alpha}^{2}t_{\beta}^{2}$, the cross term having zero mean. Summing over states gives the formula, and inserting (iii) with $\sigma = \tfrac12, 1, \tfrac12$ gives $c = 12 \cdot \tfrac12 \cdot \tfrac16 = 1$, $12 \cdot \tfrac14 = 3$, $12 \cdot \tfrac12 \cdot 2 = 12$ and $\tau = 2 \cdot \tfrac12 (1 - 6\xi)^{2} = (1 - 6\xi)^{2}$ for scalars. For fermion to fermion the amplitude of each relative helicity is built from the four numbers $B_{mn} = B(\mathcal{T}_{m}, R \cdot \mathcal{T}'_{n})$ as $(B_{11} \pm B_{22},\, B_{21} \mp B_{12})$, with squared length $\sum B_{mn}^{2} \pm 2(B_{11}B_{22} - B_{12}B_{21})$; this agrees at each sampled angle with the two values of the stored theorem's check. By (ii) the mean of $B_{11}B_{22}$ is $\tfrac15\ip{\mathcal{T}_{1}}{\mathcal{T}_{2}}\ip{\mathcal{T}'_{1}}{\mathcal{T}'_{2}} = 0$ and likewise for $B_{12}B_{21}$, so both relative helicities have the weight $\tfrac{4\pi}{5}(s^{2}/4)^{2}/s^{4}$. The script confirms the formula in all nine channels, at two pairs of couplings, by exact integration over the cosine of polynomials fixed at seven rational angles (degree at most four, the coefficients of degrees five and six being zero).

\emph{(v).} The map $\beta(a) = b \times a$ is skew, $x \cdot \beta(y) = -\beta(x) \cdot y$, and $L_{b}\mathcal{T} = \beta\mathcal{T} - \mathcal{T}\beta$, so $\tr L_{b}\mathcal{T} = 0$ for every $\mathcal{T}$. For a pair of weight $m \neq 0$, $m\,\tr\mathcal{T}_{2} = \tr L_{b}\mathcal{T}_{1} = 0$ and $-m\,\tr\mathcal{T}_{1} = \tr L_{b}\mathcal{T}_{2} = 0$, and $\ip{\mathcal{T}_{n}}{P_{k}} = \tr\mathcal{T}_{n}$ by (i). The second sentence of the clause applies this to any two-quantum element that is a pair of weight one.

\emph{(vi).} A source invariant under the rotors fixing $k$ has an invariant trace-free part, zero by (ii): it is a multiple of $P_{k}$. Clause (i) of the stored theorem gives $\mathcal{G}_{k}[\mathcal{X} - \tfrac12\id\tr\mathcal{X}] = \tfrac12 Q(k)\,\mathcal{X}$ for a source $\mathcal{X}$, and in four dimensions $\tr\mathcal{G}_{k}[h] = -\mathcal{R}_{k}[h]$, so $\mathcal{R}_{k}[\mathcal{X} - \tfrac12\id\tr\mathcal{X}] = -\tfrac12 Q(k)\tr\mathcal{X}$. Put $h = \tfrac{2\lambda}{Q(k)}\bigl[\mathcal{T}^{\circ} + \tfrac{\gamma}{3}\tr\mathcal{T}\,(P_{k} - \tfrac32\id)\bigr]$. Then $\mathcal{G}_{k}[h] = \lambda\bigl[\mathcal{T}^{\circ} + \tfrac{\gamma}{3}\tr\mathcal{T}\,P_{k}\bigr]$ and $\mathcal{R}_{k}[h] = -\lambda\gamma\tr\mathcal{T}$, and the added term is $\mu \cdot (-Q(k)P_{k}) \cdot (-\lambda\gamma\tr\mathcal{T})$; the sum is $\lambda\bigl[\mathcal{T}^{\circ} + \tfrac{\gamma}{3}(1 + 3\mu Q(k))\tr\mathcal{T}\,P_{k}\bigr]$, which is $\lambda\mathcal{T}$ exactly when $\gamma = 1/(1 + 3\mu Q(k))$. If $h$ solves the equation without source, its trace gives $-(1 + 3\mu Q(k))\,\mathcal{R}_{k}[h] = 0$, so $\mathcal{R}_{k}[h] = 0$ and $\mathcal{G}_{k}[h] = 0$: a displacement, by the stored kernel (Theorem \claimref{synthesis.weyl-flat-variations-carry-only-the-stress-trace-and-the-whole-stress-clause-is-the-trace-free-channel}(i); clause (i) of Theorem \claimref{synthesis.gravitational-pair-production-from-the-operator-source-at-first-order}). Displacements pair to zero with sources, so the interaction is $\ip{\mathcal{T}_{1}}{h[\mathcal{T}_{2}]} = \tfrac{2\lambda}{Q(k)}\bigl[\ip{\mathcal{T}_{1}^{\circ}}{\mathcal{T}_{2}^{\circ}} + \tfrac{\gamma}{3}t_{2}\,(t_{1} - \tfrac32 t_{1})\bigr]$, which is the stated $B_{\mu}$. The weights follow as in (iv) with $t_{\alpha}t_{\beta}$ replaced by $\gamma\,t_{\alpha}t_{\beta}$. The script verifies the solution, its scalar curvature and the interaction at four values of $\mu$ for scalar, fermion and vector sources.

\emph{The Reading} is argued in the statement. Its three inputs from the literature (the state-independence of the trace of the renormalised stress of a free field; the form of the stress trace in an interacting gauge theory; the size of the anomaly coefficient) are quoted from memory and are used in no step above.
